\section{The relocated gate}
\label{sec:tpc42-next-gate}

The two candidate shortcuts considered in TPC--41 now have exact
endpoints.  The scalar progression direction postcomposes through every
bounded residue-synchronous Bessel map on the reduced-residue sector, but
the actual terminal vectors
are orthogonal rather than synchronous within each residue fiber.  The
localized convolution completes exactly, but its boundary is the
negative local cell rather than a smaller error.  Neither conclusion is
heuristic.

\subsection{Changing only the alias modulus does not move the equality column}

Fix a common inherited physical row set, coefficient field, and terminal
support, and vary only the auxiliary modulus used in the alias shifts.
Attach that modulus to the notation and write
\(\mathscr T_k^L(M)\) for the output in
\eqref{eq:tpc42-left-column}.  The coefficients \(c^L\) are inherited
before the additive alias bank is introduced, and \(kM\) is the only
occurrence of \(M\) in the terminal target.

\begin{proposition}[Equality-column invariance]
\label{prop:tpc42-equality-column-invariance}
For every two admissible auxiliary triples \(M,M'\) applied to this same
physical coefficient field and support,
\begin{equation}
 \boxed{
 \mathscr T_0^L(M)=\mathscr T_0^L(M'),
 \qquad
 \mathscr T_0^R(M)=\mathscr T_0^R(M').}
 \label{eq:tpc42-equality-column-invariance}
\end{equation}
Consequently, for every nonempty finite admissible family
\(\mathfrak M\),
\begin{equation}
 \frac1{|\mathfrak M|}\sum_{M\in\mathfrak M}
 \left(
  \|\mathscr T_0^L(M)\|^2+
  \|\mathscr T_0^R(M)\|^2
 \right)
 =\|\mathscr T_0^L\|^2+
  \|\mathscr T_0^R\|^2.
 \label{eq:tpc42-equality-average-invariant}
\end{equation}
\end{proposition}

\begin{proof}
At \(k=0\),
\(n_{\alpha,j,0}=m_\alpha j+h\), independently of \(M\).  The
coefficient in \eqref{eq:tpc42-left-c} is also independent of the
auxiliary triple.  This proves the two vector identities and hence their
averaged norm identity.
\end{proof}

Thus an average of a nonnegative complete physical energy over auxiliary
triples can be useful only if it directly bounds the common equality
column or supplies information on an \(M\)-dependent defect that can be
transferred back to it.  Merely selecting the smallest member of the
average cannot make the equality column smaller.  Nonzero aliases do
depend on \(M\), so a future averaged theorem may still use them; the
proposition closes only the naive equality-selection argument.

The common-field qualification is substantive.  If a construction
reselects residue masks, active rows, or a regular subpacket separately
for each modulus, then those data are part of the \(M\)-dependence and
\cref{prop:tpc42-equality-column-invariance} does not apply.  The
all-residue identities of \cref{sec:tpc42-residue-orthogonality} avoid
such an implicit pruning.

\subsection{The exact positive continuation}

After projective separation, the direct route is the restricted
orbit--slope Bessel estimate
\eqref{eq:tpc42-orbit-slope-equality-target}.  In expanded form, its
off--diagonal is a weighted average of the affine products
\begin{equation}
 \mu(d_\alpha)\mu(d_\beta)
 \mu(m_\alpha j+h)\mu(m_\beta j+h),
 \qquad \alpha\ne\beta,
 \label{eq:tpc42-next-four-mobius}
\end{equation}
with the actual TPC--36 masks and source weights.  The two target affine
forms are nonproportional, but their slopes and coefficients grow with
\(X\).  Existing fixed-form or dense-shift theorems do not supply the
uniform weighted estimate.

There are three now sharply separated subroutes.
\begin{enumerate}[label=(\alph*)]
 \item \textbf{Direct cross-residue dispersion.}
 Prove \eqref{eq:tpc42-orbit-slope-equality-target} for the declared
 projective layers by a bilinear or Bessel inequality that uses the
 source signs, the affine determinants, and the orbit average jointly.
 This is the shortest route to the physical equality column.

 \item \textbf{Average the canonical defect itself.}
 For each admissible triple, estimate
 \begin{equation}
  \|T_G(I-\Pi_w)\mu\|^2
  \label{eq:tpc42-next-averaged-defect}
 \end{equation}
 rather than a scalar residue variance.  Positivity can then select a
 good triple for an \(M\)-dependent alias family.  By
 \cref{lem:tpc42-density-loss}, restriction from a dense modulus or
 shift average needs more than three powers of logarithmic saving unless
 the theorem is proved directly on the triple-prime family.

 \item \textbf{Recombine localization before taking the Hilbert norm.}
 The three-face completion of
 \cref{thm:tpc42-three-face-completion} cannot treat one local cell.
 A different argument could seek cancellation among a complete signed
 partition of all dyadic cells \emph{before} squaring.  Such a theorem
 must retain the physical \((\gamma,j)\) coordinates; scalar convolution
 after those coordinates have been summed is insufficient.
\end{enumerate}

The first route is the next minimal theorem.  The second is a genuine
selection problem with an explicit density threshold.  The third changes
the order of localization and Hilbert synthesis and is therefore a
different architecture, not a repair of the small-boundary argument.

\subsection{Claim boundary}

This paper proves no bound of the form
\eqref{eq:tpc42-orbit-slope-equality-target} or
\eqref{eq:tpc42-orbit-slope-trace-target}.  It does not prove
cancellation for an individual pair of growing affine forms, a
fixed-shift Chowla theorem, a diagonal-scale raw scalar progression
variance, or the physical four--M\"obius estimate
\eqref{eq:tpc42-inherited-gate}.  The countermodels rule out universal
black-box implications; they do not assert that the actual M\"obius
coefficients saturate those models.  The convolution lower bound concerns
an unmasked coefficient cell; it is not a lower bound for the signed
terminal output.  In particular, no sieve-parity breach,
Hardy--Littlewood asymptotic, or infinitude of twin primes is claimed.
